\chapter{System Analysis and Design}
\label{ch:design}

\section{System Architecture}
The \textit{FinMetrics Suite} is structured as a decoupled, multi-layer architecture. This decoupling ensures that the user interface (UI) layer is separated from the analytical modules, database cache, and API engines, facilitating code maintenance and extensions.

The architecture is divided into three primary layers:
\begin{enumerate}
    \item \textbf{Presentation Layer (Streamlit UI)}: Renders the sidebar controls (scopes, calibrations), interactive data tabs, Plotly charts, and file upload dropzones.
    \item \textbf{Analytical Kernel Layer (Python Packages)}: The core processing layer containing standard modules:
        \begin{itemize}
            \item \textbf{Sector Loader}: Maps sectors to public tickers.
            \item \textbf{Data Fetcher}: Retrieves Yahoo Finance statements and manages local cache.
            \item \textbf{Financial Parser}: Standardizes columns using the alias dictionary.
            \item \textbf{Ratio Engine}: Computes 30+ ratios and YoY growth rates.
            \item \textbf{Outlier Engine}: Identifies anomalies and runs archetype clustering.
            \item \textbf{Benchmark Core}: Coordinates baseline, pillar scoring, and risk categorization.
            \item \textbf{Tally Bridge}: Stream-processes ERP journals and balances Trial Balances.
            \item \textbf{Client Comparator}: Aligns client ratios with public baselines.
            \item \textbf{Benchmark Classifier}: Assigns health labels and structures the SWOT matrix.
            \item \textbf{AI Handler}: Manages Gemini ticker discovery and diagnostic chat sessions.
        \end{itemize}
    \item \textbf{Data Cache Layer}: Persistent storage in JSON and Parquet formats to avoid API rate limits.
\end{enumerate}

\section{System Workflow}
The application executes in a clean, stateful process flow divided into two phases:

\subsection{Phase 1: Listed Peer Group Baseline Generation}
\begin{enumerate}
    \item The user defines the benchmark scope via the sidebar: choosing a pre-mapped Sector, selecting an Industry, or inputting an abstract term for Gemini AI Search.
    \item The sidebar scans environment variables to configure the active Gemini API keys.
    \item If AI Search is selected, the \textit{AI Handler} queries Gemini to discover listed stock tickers and dynamically define their sector and market (appending exchange suffixes like \textit{.NS}).
    \item The \textit{Data Fetcher} checks local cache; if missing, it queries `yfinance` to pull quarterly and annual statements.
    \item The \textit{Financial Parser} resolves nomenclature variations using the 60-field \texttt{ALIAS\_MAP}, standardizing statements.
    \item The \textit{Ratio Engine} computes standard ratios.
    \item The \textit{Outlier Engine} detects statistical anomalies across all ratios using Tukey Fences and Z-scores.
    \item The \textit{Benchmark Core} masks the identified outlier cells, computes clean peer statistics, grades company ratios to assign percentile ranks, and ranks firms by their weighted composite health score.
\end{enumerate}

\subsection{Phase 2: Private Enterprise Client Benchmarking}
\begin{enumerate}
    \item The user uploads private data through the landing page dropzones (raw Tally ERP CSV files including Ledger Master and Journal Vouchers) or manually enters financials/ratios.
    \item If Tally files are uploaded, the \textit{Tally Bridge} stream-processes the multi-thousand line journals, maps groups to internal pillars, balances the Trial Balance, and auto-computes the client's ratios.
    \item The \textit{Client Comparator} normalizes the client's ratios and maps them to the frozen public peer distribution to determine their percentile rank.
    \item The K-Means algorithm clusters the peer set and maps the client to a corporate archetype (e.g., \textit{Elite Market Leader}).
    \item The \textit{Benchmark Classifier} maps the client's strengths and weaknesses into a 2x2 SWOT grid.
    \item The user can download a styled multi-sheet Excel report or launch an stateful diagnostic chat session with Gemini to analyze operational improvements.
\end{enumerate}

\begin{figure}[htbp]
    \centering
    \begin{tikzpicture}[
        title/.style={font=\bfseries\small, text centered},
        block/.style={rectangle, draw=black, fill=blue!5, text width=4.5cm, text centered, rounded corners, minimum height=0.9cm, font=\scriptsize, align=center},
        arrow/.style={thick, -{Stealth[length=5pt, width=3.75pt]}}
    ]
        % Top Streamlit UI (Centered at x=0, y=0)
        \node (ui) [block, fill=orange!10, text width=6cm, minimum height=1.0cm] at (0,0) {\textbf{Streamlit UI (Visual Control Deck)}};
        
        % Left Column (Phase 1: Public Peer Baseline, Centered at x=-3.5)
        \node (yfinance) [block] at (-3.5, -2.0) {\textbf{yfinance API}\\Harvest real-time public statements};
        \node (parser) [block] at (-3.5, -3.5) {\textbf{Financial Parser}\\Standardize nomenclature via 60-field alias map};
        \node (ratio) [block] at (-3.5, -5.0) {\textbf{Ratio Engine}\\Compute 30+ canonical ratios \& YoY/QoQ growth};
        \node (outlier) [block] at (-3.5, -6.5) {\textbf{Outlier Engine}\\IQR, Z-Score, and Isolation Forest detection};
        \node (core) [block, fill=blue!10] at (-3.5, -8.0) {\textbf{Benchmark Core}\\Mask outliers, build clean baselines, rank peers};
        
        % Right Column (Phase 2: Client Integration, Centered at x=3.5)
        \node (ingest) [block] at (3.5, -2.0) {\textbf{TallySync / Manual Entry}\\Ingest private ERP journals or manually enter financials};
        \node (tally) [block] at (3.5, -3.5) {\textbf{Tally Bridge}\\Reconstruct Trial Balance \& autocompute ratios};
        \node (comp) [block] at (3.5, -5.0) {\textbf{Client Comparator}\\Overlay client ratios against public distribution};
        \node (swot) [block] at (3.5, -6.5) {\textbf{SWOT \& Archetype}\\Run KMeans clustering \& classify SWOT categories};
        \node (report) [block, fill=green!10] at (3.5, -8.0) {\textbf{Report Generator}\\Export interactive visual dashboards \& styled Excel};
        
        % Connections using a clean drop-down split bus
        \coordinate (ui_split) at ($(ui.south) + (0,-0.4)$);
        \draw [thick] (ui.south) -- (ui_split);
        \draw [thick] (-3.5, -0.9) -- (3.5, -0.9);
        \draw [arrow] (-3.5, -0.9) -- (yfinance.north);
        \draw [arrow] (3.5, -0.9) -- (ingest.north);
        
        \draw [arrow] (yfinance) -- (parser);
        \draw [arrow] (parser) -- (ratio);
        \draw [arrow] (ratio) -- (outlier);
        \draw [arrow] (outlier) -- (core);
        
        \draw [arrow] (ingest) -- (tally);
        \draw [arrow] (tally) -- (comp);
        \draw [arrow] (comp) -- (swot);
        \draw [arrow] (swot) -- (report);
    \end{tikzpicture}
    \caption{Dual-Phase Platform Architecture and Workflow}
    \label{fig:workflow_flowchart}
\end{figure}

\section{Data Flow Diagrams}
To detail the data movement within the platform, DFD Level 0 (Context Diagram) and DFD Level 1 (Process Diagram) are provided below.

\subsection{DFD Level 0: Context Diagram}
The Level 0 Data Flow Diagram (Figure~\ref{fig:dfd0}) outlines the primary external entities interacting with the FinMetrics Suite.

\begin{figure}[htbp]
    \centering
    \begin{tikzpicture}[
        node distance=1.5cm and 2.2cm,
        entity/.style={rectangle, draw, fill=orange!10, rounded corners, minimum width=2.4cm, minimum height=1.2cm, font=\small\bfseries, text centered, align=center},
        process/.style={rectangle, draw, fill=blue!10, rounded corners, minimum width=3.8cm, minimum height=1.8cm, font=\small\bfseries, text centered, align=center},
        arrow/.style={thick,->,>=stealth},
        label/.style={font=\scriptsize, align=center}
    ]
        % Nodes
        \node (user) [entity] {User};
        \node (suite) [process, right=of user] {FinMetrics Suite};
        \node (yfinance) [entity, right=of suite] {yfinance\\API};
        
        % Left side arrows (User <-> Suite)
        \draw [arrow] ($(user.east) + (0, 0.25)$) -- node[above, label]{Ingest /\\Config} ($(suite.west) + (0, 0.25)$);
        \draw [arrow] ($(suite.west) - (0, 0.25)$) -- node[below, label]{PDF/Excel\\Output} ($(user.east) - (0, 0.25)$);
        
        % Right side arrows (Suite <-> yfinance)
        \draw [arrow] ($(suite.east) + (0, 0.25)$) -- node[above, label]{Ticker\\Query} ($(yfinance.west) + (0, 0.25)$);
        \draw [arrow] ($(yfinance.west) - (0, 0.25)$) -- node[below, label]{Statements} ($(suite.east) - (0, 0.25)$);
    \end{tikzpicture}
    \caption{Data Flow Diagram Level 0: Context Diagram}
    \label{fig:dfd0}
\end{figure}

\subsection{DFD Level 1: Process Diagram}
The Level 1 Data Flow Diagram (Figure~\ref{fig:dfd1}) details the internal processes, data repositories, and transitions.

\begin{figure}[htbp]
    \centering
    % Soft pastel colors defined locally inside this figure to prevent global namespace clash
    \definecolor{peachbg}{RGB}{253, 238, 226}
    \definecolor{lavenderbg}{RGB}{238, 240, 252}
    \definecolor{greenbg}{RGB}{228, 248, 232}
    \definecolor{boxborder}{RGB}{40, 45, 55}
    \resizebox{\textwidth}{!}{%
    \begin{tikzpicture}[
        % Serif font to maintain the academic/formal aesthetic
        font=\rmfamily\small,
        >={Stealth[length=6.5pt, width=4.5pt]},
        % Unified base styling for all nodes
        box/.style={
            draw=boxborder, rounded corners=5pt, align=center,
            minimum height=1.6cm, inner sep=6pt, thick
        },
        % Specific styles inheriting from the base
        lavender/.style={box, fill=lavenderbg, minimum width=2.8cm},
        peach/.style={box, fill=peachbg, minimum width=2.2cm},
        green/.style={box, fill=greenbg, minimum width=2.8cm}
    ]

        % =========================================================================
        % 1. NODES (Placed on a strict coordinate grid for perfect alignment)
        % =========================================================================
        
        % Far Left
        \node[peach] (user) at (-5.5, 0) {
            \textbf{User}
        };

        % Center Left
        \node[lavender, minimum width=3.4cm] (n1) at (0, 0) {
            \textbf{1.0}\\[1pt]
            \textbf{Data Ingest}\\[1pt]
            \textbf{\& Parsing}
        };

        % Top Center (API & Cache)
        \node[peach] (yfin) at (0, 4.0) {
            \textbf{yfinance}\\[1pt]
            \textbf{API}
        };
        \node[green] (d1) at (5.5, 4.0) {
            \textbf{D1: Raw Cache}
        };

        % Center Right & Far Right
        \node[lavender] (n2) at (5.5, 0) {
            \textbf{2.0}\\[1pt]
            \textbf{Ratio}\\[1pt]
            \textbf{Computation}
        };
        \node[lavender] (n3) at (10.5, 0) {
            \textbf{3.0}\\[1pt]
            \textbf{Outlier}\\[1pt]
            \textbf{Scrubbing}
        };

        % Bottom Center
        \node[lavender] (n4) at (5.5, -3.2) {
            \textbf{4.0}\\[1pt]
            \textbf{Scoring \&}\\[1pt]
            \textbf{Benchmarking}
        };


        % =========================================================================
        % 2. CONNECTIONS & FLOWS (Orthogonal routing)
        % =========================================================================
        \begin{scope}[thick, draw=boxborder!90]

            % Horizontal Main Track
            \draw[->] (user.east) -- node[midway, below, font=\scriptsize, draw=none, fill=none] {ERP / Manual} (n1.west);
            \draw[->] (n1.east) -- (n2.west);
            \draw[->] (n2.east) -- (n3.west);

            % Vertical API Bi-directional Flow 
            % 'sloped' forces the text to align vertically with the arrow paths.
            % 'above=2pt' adds a tiny buffer so the text doesn't touch the line itself.
            
            \draw[->] ([xshift=-15pt]n1.north) -- 
                node[midway, sloped, above=2pt, font=\scriptsize, draw=none, fill=none] {Ticker} 
                ([xshift=-15pt]yfin.south);
                
            \draw[->] ([xshift=15pt]yfin.south) -- 
                node[midway, sloped, above=2pt, font=\scriptsize, draw=none, fill=none] {Raw HTML} 
                ([xshift=15pt]n1.north);

            % Flow to D1 Cache (Routed to cleanly bypass the API box)
            % Starts at the far right edge of the top of 1.0, goes up, then right into D1
            \draw[->, rounded corners=4pt] ([xshift=45pt]n1.north) |- (d1.west);

            % D1 Cache to 2.0 (Perfect vertical drop)
            \draw[->] (d1.south) -- (n2.north);

            % 3.0 to 4.0 (Orthogonal corner drop)
            \draw[->, rounded corners=4pt] (n3.south) |- (n4.east);

            % 4.0 Return path back to User
            % Goes straight left, then straight up into the bottom of the User node
            \draw[->, rounded corners=4pt] (n4.west) -- 
                node[midway, below, font=\scriptsize, draw=none, fill=none] {SWOT / Ranks} 
                (-5.5, -3.2) -- (user.south);

        \end{scope}

    \end{tikzpicture}%
    }
    \caption{Data Flow Diagram Level 1: Process Diagram}
    \label{fig:dfd1}
\end{figure}

\section{Module Descriptions}
The system consists of the following modules, each encapsulating specific functional logic:
\begin{itemize}
    \item \textbf{sector\_loader.py}: A template mapping utility that lists standard public market tickers under 13 sectors: Automobiles, IT, FMCG, Banking, Pharma, Oil \& Gas, Metals, Cement, Telecom, Power, Realty, Capital Goods, and Infrastructure.
    \item \textbf{data\_fetcher.py}: Handles secure connections to the `yfinance` Ticker wrapper. Implements asynchronous API requests and caches raw quarterly/annual statements in Parquet format with a 12-hour TTL config.
    \item \textbf{financial\_parser.py}: Normalizes messy statements into clean, structured DataFrames. Scans the rows, renames matched variants using the alias map \texttt{ALIAS\_MAP}, removes duplicate metrics, and converts all entries to numeric types.
    \item \textbf{ratio\_engine.py}: Performs financial ratio computations. Reconstructs quarterly time-series, calculates YoY/QoQ percentage changes, and computes CAGR growth trends.
    \item \textbf{outlier\_engine.py}: The statistical processing module. Calculates Tukey fences, robust Z-scores, and runs K-Means peer clustering to map companies into corporate archetypes.
    \item \textbf{benchmark\_core.py}: Orchestrates the baseline calculation. Masks outlier cells, computes clean peer statistics (medians, deciles), and ranks firms by their weighted composite health score.
    \item \textbf{tally\_bridge.py}: The ERP double-entry bridge. Ingests ledger groupings, streams journal transaction CSVs, balances Trial Balances, and maps entries to standard public financial pillars.
    \item \textbf{client\_comparator.py}: Integrates client records by mapping private ERP journal outputs or manually entered financials/ratios to the public sector distribution and calculating precise percentiles.
    \item \textbf{benchmark\_classifier.py}: The business logic translator. Assigns health labels and structures the 2x2 SWOT grid.
    \item \textbf{report\_generator.py}: Generates styled Excel workbooks using custom colors, column auto-fitting, and formatted dashboards.
    \item \textbf{ai\_handler.py}: Integrates with the Gemini API to search for sector tickers and run the interactive diagnostic chat.
\end{itemize}
